\section{Dice Combinations}
\label{sec:cses-dice-combinations}

\problemstatement{Count the number of ways to construct the sum $n$ by
throwing a die one or more times, where each throw yields $1$ to $6$ and
different \emph{orders} count as different ways. Output the count modulo
$10^9+7$.}

\begin{patternbox}
\emph{``Number of ways'' + ``order matters''} is a one-dimensional
counting DP over the target sum. Let $\textit{dp}[x]$ be the number of
ordered sequences summing to $x$; the last throw was some value
$1\ldots6$, so $\textit{dp}[x]$ sums the ways to have reached
$x-1,\ldots,x-6$ before that final throw.
\end{patternbox}

\subsection*{State and transition}

Define $\textit{dp}[x]$ as the number of ways to obtain sum $x$. The
final die throw contributes a value $k\in\{1,\ldots,6\}$, and everything
before it must sum to $x-k$. Since every different final value (and every
different earlier sequence) is a distinct ordered way, we simply add:
\[
  \textit{dp}[x]=\sum_{k=1}^{6}\textit{dp}[x-k],\qquad \textit{dp}[0]=1.
\]
The base $\textit{dp}[0]=1$ counts the single empty sequence.

\begin{ideabox}[Ordered Counting: Sum Over the Last Move]
Because order matters, we classify sequences by their \emph{last} throw.
Summing over all six possible last throws counts every ordered sequence
exactly once, with no double counting -- each sequence has a unique final
element.
\end{ideabox}

\subsection*{Algorithm}

\begin{enumerate}
  \item $\textit{dp}[0]=1$.
  \item For $x=1\ldots n$: $\textit{dp}[x]=\sum_{k=1}^{6}
        \textit{dp}[x-k]$ (skipping $x-k<0$), taken modulo $10^9+7$.
  \item Answer is $\textit{dp}[n]$.
\end{enumerate}

\subsection*{Dry run}

\dryrun{$n=3$.}

\begin{itemize}
  \item $\textit{dp}[0]=1$.
  \item $\textit{dp}[1]=\textit{dp}[0]=1$ (just ``1'').
  \item $\textit{dp}[2]=\textit{dp}[1]+\textit{dp}[0]=2$ (``1+1'',
        ``2'').
  \item $\textit{dp}[3]=\textit{dp}[2]+\textit{dp}[1]+\textit{dp}[0]
        =2+1+1=4$ (``1+1+1'',``1+2'',``2+1'',``3'').
\end{itemize}

\subsection*{C++ implementation}

\begin{lstlisting}
#include <bits/stdc++.h>
using namespace std;
const int MOD = 1e9 + 7;

int main() {
    int n;
    cin >> n;
    vector<long long> dp(n + 1, 0);
    dp[0] = 1;
    for (int x = 1; x <= n; x++)
        for (int k = 1; k <= 6; k++)
            if (x - k >= 0) dp[x] = (dp[x] + dp[x - k]) % MOD;
    cout << dp[n] << "\n";
    return 0;
}
\end{lstlisting}

\begin{complexitybox}
\textbf{Time Complexity.} $O(6n)=O(n)$ -- each of $n$ states sums six
predecessors. \textbf{Space Complexity.} $O(n)$ for the DP array (reducible
to $O(1)$ by keeping only the last six values).
\end{complexitybox}

\begin{takeawaybox}
The archetypal ordered-counting DP: \emph{sum over the last move}. Every
``count sequences where order matters'' problem in this chapter is a
variation on this single recurrence -- the only thing that changes is the
set of allowed last moves.
\end{takeawaybox}
